\section{Morse Theory}
For everything concerning Morse theory our main reference will be the wonderful book \cite{banyaga2004lectures}. 
In Morse Theory, we will always deal with asmooth funkition $f:M\to \R$ from a smooth Riemannian manifold $(M,g)$ that is also compact. We will denote its dimension by $m=\dim (M)$.
\begin{definition}[Kritical Points]
	We call a point $p\in M$ \textbf{critical}, if $\restr{\diff f}{p}$ vanishes meaning that 
	\begin{equation*}
		\restr{\diff f}{p} :T_pM\to T_{f(p)}\R
	\end{equation*} is trivial. By $\Crit (f)$ we denote the \textbf{set of critical points}.
\end{definition}
\begin{definition}[The Hessian]
Let $f:M\to \R$ be a smooth funktion. The \textbf{Hessian} of $f$ at $p\in \Crit(f)$ is a symmetric bilinear funktion
\begin{equation*}
	\Hess_p(f):T_pM\times T_pM \to \R
\end{equation*}  that is defined as follows. Let $\xi_p, \zeta_p\in T_pM$ be two tangend vectors. Now choos two vectorfields $\Xi$ and $Z$ in a neighbourhoof of $p$ such that $\Xi(p)=\xi_p$ and $Z(p)=\zeta_p$. With this define the Hessian to be: 
\begin{equation*}
	\Hess_p(f)(\xi_p,\zeta_p)=\left(\Xi\cdot (Z\cdot f)\right)(p)
\end{equation*} Compare \ref{cor: Vectorfields hav funktions as inputs} for the notation. 
\end{definition}
\begin{definition}
	A smooth function $f: M\to \R$ is called \textbf{Morse}, if for any critical point.
\end{definition}
The Hessian is a Symmetric bilinear and non-degnereate Map. For such a map there exists the following theorem: 
\begin{theorem}[Sylvesters Law of Inertia]
	Let $A\in \Mat(n,\R)$ be symmetric and let $T,T'\in \GL(n,\R)$ and $k,k',l,l'\in \N$ such that \\
	\begin{align*}
		T^t\circ A\circ T=	\begin{pmatrix}
			1_k& 0 &0\\
			0& -1_l&0\\
			0&0&0_{n-k-l}
		\end{pmatrix}\text{ and } 
		T'^T\circ A \circ T'=
		\begin{pmatrix}
			1_{k'}& 0 &0\\
			0& -1_{l'}&0\\
			0&0&0_{n-k'-l'}
		\end{pmatrix}
	\end{align*}
	then $k=k'$, $l=l'$ and $\rank(A)=k+l.$
\end{theorem}
\begin{proof}
	Since $T,T'$ are invertible we have that
	\begin{align}
		k+l=\rank(T^t\circ A\circ T)=\rank(A)=\rank(T'^T\circ A\circ T')=k'+l' \, .
	\end{align} Hence it suffices to show that $k=k'$. Which we will do by proofing the claim:
	\begin{equation*}
		k=\max \left\{ \dim(U)|U\subseteq \R^n \text{subspace such that } x^tAx>0 ~\forall x\in U\setminus \{0\}  \right\}
	\end{equation*}
	So we start by showing \glqq $\leq $\grqq:
	Denote the first $k$ colomus of $T$ with $x_1,...,x_k$. They form a basis of $\R^n$ and with $ 0 \neq x=\sum_{i=1}^k\lambda ix^i$ we have that by bilinearity
	\begin{equation}
		x^TAx=\sum_{i=1}^{k}\lambda_i x_i^T A x=\sum_{i,j=1}^{k}\lambda_i \lambda_j x_i^T A x_j=\sum_{i=1}^{k}(\lambda_i)^2\geq 0 \,. 
	\end{equation} This conculdes the first inequality. 
	
	Now let $U$ be any $k-$dimensional subspace such that for all non zero $x\in U$ $x^TAx>0$. By a calculation analog to the one above we have that for any $x\in W:= \gen (x_{k+1},\dots x_n)$ the number $x^tAx$ is less or equal to zero. Hence, $W\cap U= \{0\}$ and with this we conclude:
	\begin{equation}
		\dim(U)=\dim(U+W)-\dim(W)+\dim(U\cap W)\leq (k+(n-k))-(n-k)+0=k\, .
	\end{equation} This is the last inequality proving the statement.
\end{proof}
